\documentclass[10pt,twocolumn]{article}
\usepackage[margin=0.75in]{geometry}
\usepackage{booktabs}
\usepackage{array}
\usepackage{amsmath}
\usepackage[hidelinks]{hyperref}
\usepackage{xcolor}
\usepackage{titlesec}
\titlespacing*{\section}{0pt}{1.2ex}{0.6ex}
\titlespacing*{\subsection}{0pt}{1.0ex}{0.4ex}
\setlength{\parskip}{0.35ex}
\newcommand{\code}[1]{\texttt{\small #1}}

\title{\vspace{-3ex}\textbf{Which LLM Should Score Your CVEs?}\\
\large A Cost-Aware Benchmark for Context-Conditioned\\ Vulnerability Prioritization}
\author{Fahed Dorgaa\\\small venslabs \;\textbar\; \code{github.com/venslabs/vens-benchmark}}
\date{}

\begin{document}
\maketitle

\begin{abstract}
\noindent
A scanner hands you 400 CVEs and their CVSS base scores. Most of them do not
matter for \emph{your} deployment, and CVSS was never meant to tell you which.
\code{vens} closes that gap by asking an LLM to re-score each finding against a
short project-context file (exposure, data sensitivity, criticality, controls,
compliance) and the scanner report, producing an OWASP Risk Rating from 0--81.
The obvious question from operators is: \emph{which model do I put behind it, and
is it worth the API bill?} We benchmark eight 2026 cloud models and four local
models on the two things \code{vens} actually needs---reading a CVE and using the
context---and compare every one against a non-LLM rule engine it has to beat.
The headline: accuracy is a solved, cheap problem (a \$0.49 model ties a \$4.12
one), context use is real but uneven (reachability and accountability move
correctly; a denial-of-service cue is mostly keyword-matched), and the weakest
models---including every local one we tried---are statistically
indistinguishable from a constant guess. We give a concrete model-selection
table and release the harness so a practitioner can rank any new model in an
afternoon.
\end{abstract}

\section{Introduction}
Vulnerability scanners are good at finding CVEs and bad at telling you which to
fix first. CVSS base scores, the usual sort key, are a property of the
vulnerability, not of the system it sits in: the same 9.8 RCE is a fire drill on
an internet-facing payment API and a non-event in an air-gapped batch job. Teams
end up hand-triaging, or worse, patching by score and running out of time.

\code{vens}\footnote{\url{https://github.com/venslabs/vens}} is an
open-source tool that re-prioritizes a scan with an LLM. It takes two inputs
only---the scanner report and a hand-written \code{config.yaml} describing the
project---and emits a CycloneDX VEX document with an OWASP Risk Rating per CVE.
It has no runtime access; the context file is the operator's declared truth and
the CVE description is read to connect the two. Because the scoring is one LLM
call per finding, the choice of model is a direct cost and quality knob that
operators have to set with no data to go on.

This paper is that data. We ask:
\begin{itemize}\itemsep0.1ex
\item[\textbf{RQ1}] How well do current models read a CVE (predict its severity
  from the text alone), and how cheaply?
\item[\textbf{RQ2}] Do they actually \emph{use} the project context, or just
  echo the scanner?
\item[\textbf{RQ3}] When a model gets a context case right, is it reasoning or
  keyword-matching the CVE text?
\item[\textbf{RQ4}] What should a practitioner run, given a budget?
\item[\textbf{RQ5}] Are free local models a viable option?
\end{itemize}

\section{Background: what vens asks the model}
\label{sec:bg}
\code{vens} scores four OWASP Risk Rating factors, each 0--9, and multiplies
likelihood by impact:
\[
\text{Risk} = \underbrace{\tfrac{\text{ThreatAgent}+\text{Vulnerability}}{2}}_{\text{likelihood}}
\times
\underbrace{\tfrac{\text{Technical}+\text{Business}}{2}}_{\text{impact}} \in [0,81].
\]
Three of the four factors are close to arithmetic on the context file:
\emph{ThreatAgent} tracks the exposure band, \emph{Technical} tracks
data sensitivity, \emph{Business} tracks criticality plus a compliance bump. Only
\emph{Vulnerability}---how exploitable this CVE is in this deployment---needs the
model to read and reason. This is exactly why a non-LLM baseline is not a
strawman (\S\ref{sec:design}): most of the score can be produced by a lookup
table, so the model has to earn its cost on the parts that cannot.

\section{Related work}
CTIBench~\cite{ctibench} is a CTI benchmark suite whose CTI-VSP task asks a model
to produce a CVSS v3.1 vector from a CVE description; we reuse it directly for
RQ1 so our accuracy numbers sit next to its published baselines. Databricks'
VulnWatch~\cite{vulnwatch} prioritizes CVEs with an LLM but scores the
vulnerability in isolation; EPSS~\cite{epss} predicts exploitation probability
from global signals, again context-free. Reachability tools
(callgraph/VEX generators) answer ``can this code even run?'' but not ``does it
matter here?''. \code{vens}'s contribution, and what no public dataset labels, is
risk \emph{conditioned on a declared deployment context}; RQ2--RQ3 are, to our
knowledge, the first benchmark of that specific behavior.

\section{Benchmark design}
\label{sec:design}
We split ``best model for \code{vens}'' into the two skills it composes.

\paragraph{Half A --- CVE understanding (RQ1).} The CTI-VSP task: 200 CVEs,
predict the CVSS v3.1 vector from the description. We parse the vector with the
\code{cvss} library and score the base value, so the model is never asked to do
arithmetic. Metric is mean absolute deviation (MAD) from the NVD base score.

\paragraph{Half B --- context use (RQ2--RQ3).} 35 scenarios, each a one-CVE
Trivy report plus a \code{config.yaml}, run through the real \code{vens generate}
CLI with attestation on so we capture the four factors and the reasoning. Most
scenarios are \emph{contrastive}: hold everything fixed, vary one context axis,
and check the score moves the right way (exposure ladder, sensitivity ladder,
controls on/off, reachability declared in operator notes). Two are real showcase
CVEs. Five are \emph{discriminators}---cases where the naive lookup rule is
wrong, so a model only gets them right by reading the CVE: a pure
denial-of-service against a critical-data service (availability, not
confidentiality), an accountability bug under low data sensitivity, and an
irrelevant WAF in front of a local-only exploit.

\paragraph{The construct-validity trap (RQ3).} A discriminator can be passed by
keyword-matching (``the description says \emph{denial of service}'') rather than
reasoning. For the DoS, accountability, and irrelevant-control families we add a
\emph{non-telegraphed twin}: same context, but the CVE text no longer names the
impact class or asserts ``no confidentiality impact.'' If a model's factor
survives de-telegraphing, it was reasoning; if it drifts back toward the naive
value, it was pattern-matching.

\paragraph{The baseline it must beat.} A non-LLM rule engine reproduces
\code{vens}'s own prompt arithmetic (Technical=sensitivity, Business=criticality
$+$compliance, ThreatAgent=exposure band) and \emph{copies the vendor severity}
for the Vulnerability factor---because a rule engine cannot judge exploitability.
This is a null model, not an oracle: divergence from it is exactly where the LLM
adds value. We also report a constant ``always 7.5'' predictor as an accuracy
floor.

\paragraph{Protocol.} Temperature 0, seed 7, three repeats per cell. Models:
OpenAI \code{gpt-5.4-nano/mini}, \code{gpt-5.5}; Anthropic
\code{claude-haiku-4-5}, \code{claude-sonnet-4-6}; Google
\code{gemini-2.5-flash-lite/flash/pro}; and four local models via
Ollama\footnote{\url{https://ollama.com}} (\code{llama3.2}, \code{qwen2.5:7b},
\code{gemma2:9b}, \code{deepseek-r1:8b}). 95\% CIs are bootstrap over the 200
CVEs where per-CVE predictions were saved.

\section{Results}

\subsection{RQ1 --- reading a CVE is a solved, cheap problem}
Table~\ref{tab:mad} ranks all twelve models. The four strongest cloud models
(\code{sonnet}, \code{gpt-5.5}, \code{gpt-5.4-mini}, \code{gemini-2.5-pro}) sit
within overlapping confidence intervals---this is a \emph{tier}, not a ranking.
Two facts matter for operators. First, \code{gpt-5.4-mini} (\$0.48/200) ties
\code{gpt-5.5} (\$4.12/200): on the Pareto front of accuracy $\times$ cost,
\textbf{\code{gpt-5.5} is dominated}---\code{sonnet} is both more accurate and
cheaper. Second, every 2026 model beats 2024's GPT-4 (MAD 1.31), but the cheap
tier still trails 2024's Gemini-1.5 (1.09); progress is real but not uniform.
Accurate models slightly under-rate (bias $-0.15$); the cheap Gemini models
over-rate (up to $+0.82$), which for a triage tool is the safer direction but
inflates the queue.

\begin{table}[t]\centering\small
\caption{Half A: CVSS prediction on 200 CVEs. MAD lower is better; a constant
7.5 guess scores 1.573.}
\label{tab:mad}
\begin{tabular}{@{}lrrr@{}}
\toprule
Model & MAD & 95\% CI & \$/200 \\
\midrule
claude-sonnet-4-6      & \textbf{0.746} & [0.62, 0.88] & 1.88 \\
gpt-5.5                & 0.751 & --- & 4.12 \\
gpt-5.4-mini           & 0.836 & [0.69, 1.00] & \textbf{0.48} \\
gemini-2.5-pro         & 0.879 & [0.73, 1.03] & 1.16 \\
gemini-2.5-flash-lite  & 1.149 & [0.98, 1.32] & 0.07 \\
gemini-2.5-flash       & 1.168 & [1.00, 1.33] & 0.34 \\
gpt-5.4-nano           & 1.195 & --- & 0.21 \\
claude-haiku-4-5       & 1.288 & --- & 0.52 \\
\midrule
\multicolumn{4}{@{}l}{\emph{local (Ollama, free):}}\\
deepseek-r1:8b         & 1.479 & [1.26, 1.72] & 0.00 \\
qwen2.5:7b             & 1.536 & [1.35, 1.74] & 0.00 \\
gemma2:9b              & 1.643 & [1.45, 1.84] & 0.00 \\
llama3.2               & 1.686 & [1.53, 1.85] & 0.00 \\
\midrule
\emph{constant 7.5 (floor)} & 1.573 & --- & 0.00 \\
\bottomrule
\end{tabular}
\end{table}

\subsection{RQ2 --- context use is real but uneven}
On the two showcases the good models do the hard thing. A Log4Shell finding
(CVE-2021-44228, CVSS 10.0) in a low-value, internal, message-lookup-disabled
context drops to LOW on 7 of 8 cloud models; a medium axios token leak
(CVE-2023-45857) on an internet PII portal under GDPR rises to HIGH on all 8.
The one failure is telling: \code{gemini-2.5-flash-lite} \emph{re-inflates}
Log4Shell to 48.8/Critical---worse than the rule baseline, which de-escalates it
correctly to LOW. Cheapness has a price.

The contrastive ladders (exposure, sensitivity, criticality, compliance) are
monotone for nearly every model. The discriminators separate the field:
\emph{reachability} declared in operator notes moves the score for 6/8 models
(\code{flash-lite} and \code{nano} ignore the notes entirely, scoring reachable
and unreachable identically); \emph{accountability} lifts the Technical factor on
8/8; the \emph{irrelevant WAF} is correctly \emph{not} credited against a
local-only exploit by all. Per-model, a Wilcoxon test of (LLM $-$ baseline) risk
is significant for \code{gemini-2.5-flash}, \code{-pro}, and \code{gpt-5.4-mini};
crucially, \code{flash-lite} and \code{nano} have a median divergence of
\emph{zero}---on context, they add nothing over the lookup table.

\subsection{RQ3 --- some ``context use'' is keyword-matching}
Table~\ref{tab:ct} is the honest part. When we strip the giveaway phrase, the
\textbf{accountability} signal holds---Technical stays at 8--9 (naive rule: 2)
with near-zero drift, so models genuinely infer that tampering with an audit
trail is an integrity loss. The \textbf{DoS} signal does not: Technical was only
marginally below the naive 9 to begin with ($\sim$8), and it drifts back up when
``denial of service'' is removed. Under three different baseline threshold sets
the accountability conclusion survives and the DoS one does not. So ``the model
understands availability-only impact'' is largely an artifact of the CVE saying
so; report it as weak.

\begin{table}[t]\centering\small
\caption{RQ3: mean Technical factor, telegraphed vs.\ non-telegraphed CVE
(naive rule in parentheses). Small drift = reasoning; upward drift = keyword-matching.}
\label{tab:ct}
\begin{tabular}{@{}lcc@{}}
\toprule
 & Accountability (2) & DoS (9) \\
Model & tel.\,$\to$\,non-tel. & tel.\,$\to$\,non-tel. \\
\midrule
claude-sonnet-4-6 & 9.0 $\to$ 9.0 & 8.0 $\to$ 8.0 \\
gpt-5.5           & 8.0 $\to$ 8.0 & 8.0 $\to$ 8.0 \\
gpt-5.4-mini      & 8.3 $\to$ 8.0 & 7.7 $\to$ 8.0 \\
gemini-2.5-pro    & 8.0 $\to$ 8.0 & 8.0 $\to$ 8.0 \\
claude-haiku-4-5  & 8.3 $\to$ 8.0 & 7.0 $\to$ 8.5 \\
\midrule
\emph{verdict} & \textbf{robust} & \textbf{fragile} \\
\bottomrule
\end{tabular}
\end{table}

\subsection{RQ4 --- what to run}
Table~\ref{tab:reco} is the practitioner takeaway. Consistency is the sample
standard deviation of the risk over three repeats; it is a reliability floor, not
a quality measure (a model can be reproducibly wrong), but for auditable VEX you
want it low. \code{claude-sonnet-4-6} is the default: best accuracy, second-most
stable (SD 0.85), on the Pareto front. \code{gpt-5.4-mini} is the value pick but
jittery (SD 3.19)---run $n\ge3$ and take the median. \code{gemini-2.5-flash-lite}
is the throwaway-triage option at \$0.07/200, with the caveat that it re-inflates
the hard case and adds little on context. \code{gpt-5.5} and \code{gpt-5.4-nano}
have no niche: one is dominated, the other is jittery \emph{and} adds nothing.

\begin{table}[t]\centering\small
\caption{RQ4: recommendation. SD = risk stdev over 3 repeats (lower = more auditable).}
\label{tab:reco}
\begin{tabular}{@{}llc@{}}
\toprule
Use case & Model & SD \\
\midrule
Signed / audited VEX & claude-sonnet-4-6 & 0.85 \\
Best value           & gpt-5.4-mini      & 3.19 \\
Throwaway triage     & gemini-2.5-flash-lite & 0.00 \\
Avoid                & gpt-5.5, gpt-5.4-nano & --- \\
\bottomrule
\end{tabular}
\end{table}

\subsection{RQ5 --- local models are not there yet}
All four local models land below the cheapest cloud model on accuracy, and their
CIs \emph{overlap the constant-guess floor}: statistically they are
indistinguishable from predicting 7.5 for everything, and they over-rate by
$+0.8$ to $+1.4$. On context they are worse than the numbers suggest. The three
non-reasoning models are extremely stable (SD 0--0.15), but by rigidity, not
skill: \code{llama3.2} keeps Log4Shell at Critical regardless of context and
returns an identical maxed score across the whole exposure ladder. Only the
reasoning model, \code{deepseek-r1:8b}, actually engages the context---it is the
one local model that de-escalates the showcase (7.5/LOW) and reads
reachability---but it is jittery (SD 2.25) and still poor on Half A. For
context-conditioned scoring today, ``run a small local model'' is close to ``do
not use an LLM.''

\section{Discussion}
The split result is the useful one. Picking a model for CVE \emph{accuracy} is
nearly free: the cheap tier is good enough and the expensive tier buys nothing,
so spend the budget elsewhere. Picking a model for \emph{context} is where the
money goes---and where the cheap models quietly fail, not by erroring but by
collapsing onto the lookup table. An operator who benchmarks only accuracy (the
easy, public number) will over-value \code{flash-lite} and \code{nano} and get a
tool that ignores the very context that justified using an LLM. The two-half
design exists to catch exactly that.

\section{Threats to validity}
\textbf{Construct.} Half B scenarios are synthetic (fictional CVE IDs) to avoid
recall; the two showcases use real CVEs, where a model recalling the true
severity makes the de-escalation harder, not easier. RQ3's non-telegraphed twins
address the keyword-matching threat directly for two families but not
reachability (which is legitimately operator-declared) or the ladders.
\textbf{Statistics.} Pooled sign tests over model$\times$CVE cells are
pseudoreplicated; we report per-model counts instead, where $N{=}8$ makes 8/8 and
7/8 meaningful. CIs come from a single run's per-CVE errors and capture sampling
over CVEs, not API nondeterminism. \textbf{Scope.} MAD against NVD treats the NVD
vector as ground truth, which is itself context-free; that is the point of Half
B. The baseline is a faithful mirror of \code{vens}'s prompt, so ``beats
baseline'' means ``the LLM adds something over \code{vens}'s own hard-coded
arithmetic,'' not ``beats a strawman.'' \textbf{Data.} CTIBench is CC BY-NC-SA
4.0 and is fetched at runtime, not redistributed.

\section{Conclusion}
For \code{vens}, and for context-aware CVE scoring generally: buy accuracy cheap,
pay for context, and measure context directly---because the models that look
fine on the public accuracy number are the ones that ignore your config. Use
\code{claude-sonnet-4-6} if you sign the output, \code{gpt-5.4-mini} if you do
not, and skip the flagship. The harness, scenarios, and a one-line hook to
benchmark any new model are released so the next model can be placed the day it
ships.

\begin{thebibliography}{9}\small
\bibitem{ctibench} CTIBench: A Benchmark for Evaluating LLMs in Cyber Threat
Intelligence. NeurIPS 2024. \url{https://github.com/xashru/cti-bench}
\bibitem{vulnwatch} Databricks. VulnWatch: AI-Enhanced Prioritization of
Vulnerabilities. 2024.
\bibitem{epss} J. Jacobs et al. Exploit Prediction Scoring System (EPSS). FIRST.
\url{https://www.first.org/epss/}
\end{thebibliography}

\end{document}
